% Options for packages loaded elsewhere
% Options for packages loaded elsewhere
\PassOptionsToPackage{unicode}{hyperref}
\PassOptionsToPackage{hyphens}{url}
\PassOptionsToPackage{dvipsnames,svgnames,x11names}{xcolor}
%
\documentclass[
  letterpaper,
  DIV=11,
  numbers=noendperiod]{scrreprt}
\usepackage{xcolor}
\usepackage{amsmath,amssymb}
\setcounter{secnumdepth}{5}
\usepackage{iftex}
\ifPDFTeX
  \usepackage[T1]{fontenc}
  \usepackage[utf8]{inputenc}
  \usepackage{textcomp} % provide euro and other symbols
\else % if luatex or xetex
  \usepackage{unicode-math} % this also loads fontspec
  \defaultfontfeatures{Scale=MatchLowercase}
  \defaultfontfeatures[\rmfamily]{Ligatures=TeX,Scale=1}
\fi
\usepackage{lmodern}
\ifPDFTeX\else
  % xetex/luatex font selection
\fi
% Use upquote if available, for straight quotes in verbatim environments
\IfFileExists{upquote.sty}{\usepackage{upquote}}{}
\IfFileExists{microtype.sty}{% use microtype if available
  \usepackage[]{microtype}
  \UseMicrotypeSet[protrusion]{basicmath} % disable protrusion for tt fonts
}{}
\makeatletter
\@ifundefined{KOMAClassName}{% if non-KOMA class
  \IfFileExists{parskip.sty}{%
    \usepackage{parskip}
  }{% else
    \setlength{\parindent}{0pt}
    \setlength{\parskip}{6pt plus 2pt minus 1pt}}
}{% if KOMA class
  \KOMAoptions{parskip=half}}
\makeatother
% Make \paragraph and \subparagraph free-standing
\makeatletter
\ifx\paragraph\undefined\else
  \let\oldparagraph\paragraph
  \renewcommand{\paragraph}{
    \@ifstar
      \xxxParagraphStar
      \xxxParagraphNoStar
  }
  \newcommand{\xxxParagraphStar}[1]{\oldparagraph*{#1}\mbox{}}
  \newcommand{\xxxParagraphNoStar}[1]{\oldparagraph{#1}\mbox{}}
\fi
\ifx\subparagraph\undefined\else
  \let\oldsubparagraph\subparagraph
  \renewcommand{\subparagraph}{
    \@ifstar
      \xxxSubParagraphStar
      \xxxSubParagraphNoStar
  }
  \newcommand{\xxxSubParagraphStar}[1]{\oldsubparagraph*{#1}\mbox{}}
  \newcommand{\xxxSubParagraphNoStar}[1]{\oldsubparagraph{#1}\mbox{}}
\fi
\makeatother

\usepackage{color}
\usepackage{fancyvrb}
\newcommand{\VerbBar}{|}
\newcommand{\VERB}{\Verb[commandchars=\\\{\}]}
\DefineVerbatimEnvironment{Highlighting}{Verbatim}{commandchars=\\\{\}}
% Add ',fontsize=\small' for more characters per line
\usepackage{framed}
\definecolor{shadecolor}{RGB}{241,243,245}
\newenvironment{Shaded}{\begin{snugshade}}{\end{snugshade}}
\newcommand{\AlertTok}[1]{\textcolor[rgb]{0.68,0.00,0.00}{#1}}
\newcommand{\AnnotationTok}[1]{\textcolor[rgb]{0.37,0.37,0.37}{#1}}
\newcommand{\AttributeTok}[1]{\textcolor[rgb]{0.40,0.45,0.13}{#1}}
\newcommand{\BaseNTok}[1]{\textcolor[rgb]{0.68,0.00,0.00}{#1}}
\newcommand{\BuiltInTok}[1]{\textcolor[rgb]{0.00,0.23,0.31}{#1}}
\newcommand{\CharTok}[1]{\textcolor[rgb]{0.13,0.47,0.30}{#1}}
\newcommand{\CommentTok}[1]{\textcolor[rgb]{0.37,0.37,0.37}{#1}}
\newcommand{\CommentVarTok}[1]{\textcolor[rgb]{0.37,0.37,0.37}{\textit{#1}}}
\newcommand{\ConstantTok}[1]{\textcolor[rgb]{0.56,0.35,0.01}{#1}}
\newcommand{\ControlFlowTok}[1]{\textcolor[rgb]{0.00,0.23,0.31}{\textbf{#1}}}
\newcommand{\DataTypeTok}[1]{\textcolor[rgb]{0.68,0.00,0.00}{#1}}
\newcommand{\DecValTok}[1]{\textcolor[rgb]{0.68,0.00,0.00}{#1}}
\newcommand{\DocumentationTok}[1]{\textcolor[rgb]{0.37,0.37,0.37}{\textit{#1}}}
\newcommand{\ErrorTok}[1]{\textcolor[rgb]{0.68,0.00,0.00}{#1}}
\newcommand{\ExtensionTok}[1]{\textcolor[rgb]{0.00,0.23,0.31}{#1}}
\newcommand{\FloatTok}[1]{\textcolor[rgb]{0.68,0.00,0.00}{#1}}
\newcommand{\FunctionTok}[1]{\textcolor[rgb]{0.28,0.35,0.67}{#1}}
\newcommand{\ImportTok}[1]{\textcolor[rgb]{0.00,0.46,0.62}{#1}}
\newcommand{\InformationTok}[1]{\textcolor[rgb]{0.37,0.37,0.37}{#1}}
\newcommand{\KeywordTok}[1]{\textcolor[rgb]{0.00,0.23,0.31}{\textbf{#1}}}
\newcommand{\NormalTok}[1]{\textcolor[rgb]{0.00,0.23,0.31}{#1}}
\newcommand{\OperatorTok}[1]{\textcolor[rgb]{0.37,0.37,0.37}{#1}}
\newcommand{\OtherTok}[1]{\textcolor[rgb]{0.00,0.23,0.31}{#1}}
\newcommand{\PreprocessorTok}[1]{\textcolor[rgb]{0.68,0.00,0.00}{#1}}
\newcommand{\RegionMarkerTok}[1]{\textcolor[rgb]{0.00,0.23,0.31}{#1}}
\newcommand{\SpecialCharTok}[1]{\textcolor[rgb]{0.37,0.37,0.37}{#1}}
\newcommand{\SpecialStringTok}[1]{\textcolor[rgb]{0.13,0.47,0.30}{#1}}
\newcommand{\StringTok}[1]{\textcolor[rgb]{0.13,0.47,0.30}{#1}}
\newcommand{\VariableTok}[1]{\textcolor[rgb]{0.07,0.07,0.07}{#1}}
\newcommand{\VerbatimStringTok}[1]{\textcolor[rgb]{0.13,0.47,0.30}{#1}}
\newcommand{\WarningTok}[1]{\textcolor[rgb]{0.37,0.37,0.37}{\textit{#1}}}

\usepackage{longtable,booktabs,array}
\usepackage{calc} % for calculating minipage widths
% Correct order of tables after \paragraph or \subparagraph
\usepackage{etoolbox}
\makeatletter
\patchcmd\longtable{\par}{\if@noskipsec\mbox{}\fi\par}{}{}
\makeatother
% Allow footnotes in longtable head/foot
\IfFileExists{footnotehyper.sty}{\usepackage{footnotehyper}}{\usepackage{footnote}}
\makesavenoteenv{longtable}
\usepackage{graphicx}
\makeatletter
\newsavebox\pandoc@box
\newcommand*\pandocbounded[1]{% scales image to fit in text height/width
  \sbox\pandoc@box{#1}%
  \Gscale@div\@tempa{\textheight}{\dimexpr\ht\pandoc@box+\dp\pandoc@box\relax}%
  \Gscale@div\@tempb{\linewidth}{\wd\pandoc@box}%
  \ifdim\@tempb\p@<\@tempa\p@\let\@tempa\@tempb\fi% select the smaller of both
  \ifdim\@tempa\p@<\p@\scalebox{\@tempa}{\usebox\pandoc@box}%
  \else\usebox{\pandoc@box}%
  \fi%
}
% Set default figure placement to htbp
\def\fps@figure{htbp}
\makeatother





\setlength{\emergencystretch}{3em} % prevent overfull lines

\providecommand{\tightlist}{%
  \setlength{\itemsep}{0pt}\setlength{\parskip}{0pt}}



 


\KOMAoption{captions}{tableheading}
\makeatletter
\@ifpackageloaded{bookmark}{}{\usepackage{bookmark}}
\makeatother
\makeatletter
\@ifpackageloaded{caption}{}{\usepackage{caption}}
\AtBeginDocument{%
\ifdefined\contentsname
  \renewcommand*\contentsname{Table of contents}
\else
  \newcommand\contentsname{Table of contents}
\fi
\ifdefined\listfigurename
  \renewcommand*\listfigurename{List of Figures}
\else
  \newcommand\listfigurename{List of Figures}
\fi
\ifdefined\listtablename
  \renewcommand*\listtablename{List of Tables}
\else
  \newcommand\listtablename{List of Tables}
\fi
\ifdefined\figurename
  \renewcommand*\figurename{Figure}
\else
  \newcommand\figurename{Figure}
\fi
\ifdefined\tablename
  \renewcommand*\tablename{Table}
\else
  \newcommand\tablename{Table}
\fi
}
\@ifpackageloaded{float}{}{\usepackage{float}}
\floatstyle{ruled}
\@ifundefined{c@chapter}{\newfloat{codelisting}{h}{lop}}{\newfloat{codelisting}{h}{lop}[chapter]}
\floatname{codelisting}{Listing}
\newcommand*\listoflistings{\listof{codelisting}{List of Listings}}
\makeatother
\makeatletter
\makeatother
\makeatletter
\@ifpackageloaded{caption}{}{\usepackage{caption}}
\@ifpackageloaded{subcaption}{}{\usepackage{subcaption}}
\makeatother
\usepackage{bookmark}
\IfFileExists{xurl.sty}{\usepackage{xurl}}{} % add URL line breaks if available
\urlstyle{same}
\hypersetup{
  pdftitle={remote sensing},
  pdfauthor={Norah Jones},
  colorlinks=true,
  linkcolor={blue},
  filecolor={Maroon},
  citecolor={Blue},
  urlcolor={Blue},
  pdfcreator={LaTeX via pandoc}}


\title{remote sensing}
\author{Norah Jones}
\date{2026-01-20}
\begin{document}
\maketitle

\renewcommand*\contentsname{Table of contents}
{
\hypersetup{linkcolor=}
\setcounter{tocdepth}{2}
\tableofcontents
}

\bookmarksetup{startatroot}

\chapter*{Preface}\label{preface}
\addcontentsline{toc}{chapter}{Preface}

\markboth{Preface}{Preface}

This is a Quarto book.

To learn more about Quarto books visit
\url{https://quarto.org/docs/books}.

\begin{Shaded}
\begin{Highlighting}[]
\DecValTok{1} \SpecialCharTok{+} \DecValTok{1}
\end{Highlighting}
\end{Shaded}

\begin{verbatim}
[1] 2
\end{verbatim}

\bookmarksetup{startatroot}

\chapter{}\label{section}

\#WEEK 1 : Introduction to Remote Sensing \& Spectral Analysis using
SNAP

\#\#SUMMARY Dalam praktikum ini, citra Sentinel-2 Level-2A dianalisis
menggunakan perangkat lunak SNAP untuk memahami konsep dasar
penginderaan jauh, khususnya terkait karakteristik spektral berbagai
jenis penutup lahan. Proses yang dilakukan meliputi visualisasi citra
RGB, pembuatan subset wilayah studi, serta pengambilan sampel
menggunakan polygon untuk kelas air, vegetasi, dan area urban.

Hasil analisis menunjukkan bahwa setiap jenis penutup lahan memiliki
pola reflektansi yang berbeda terhadap panjang gelombang
elektromagnetik. Perbedaan ini memungkinkan identifikasi dan klasifikasi
objek di permukaan bumi berdasarkan respon spektralnya.

\#\#ANALYSIS (SPECTRAL SIGNATURE) Analisis spektral dilakukan
menggunakan Spectrum View untuk membandingkan reflektansi air, vegetasi,
dan area urban. Air menunjukkan nilai reflektansi yang sangat rendah
pada seluruh panjang gelombang dan kurva yang relatif datar. Hal ini
disebabkan oleh kemampuan air dalam menyerap sebagian besar radiasi
elektromagnetik, terutama pada wilayah near-infrared (NIR) dan shortwave
infrared (SWIR). Vegetasi menunjukkan pola yang berbeda, dengan
reflektansi rendah pada wilayah visible akibat penyerapan oleh klorofil,
diikuti oleh peningkatan tajam pada sekitar 700--800 nm (red edge) dan
nilai tinggi pada NIR. Setelah itu, reflektansi menurun pada SWIR akibat
pengaruh kandungan air dalam jaringan tanaman. Sebaliknya, area urban
menunjukkan nilai reflektansi sedang dengan pola yang lebih bervariasi.
Hal ini mencerminkan sifat heterogen area perkotaan yang terdiri dari
berbagai material seperti beton, aspal, dan atap bangunan. Perbandingan
ini menunjukkan bahwa vegetasi merupakan kelas yang paling mudah
diidentifikasi karena memiliki lonjakan reflektansi yang jelas pada NIR,
sementara urban merupakan yang paling kompleks karena adanya variasi
material. \#\#SCATTER PLOT ANALYSIS Scatter plot antara band merah (B4)
dan near-infrared (B8) menunjukkan adanya pemisahan kelas penutup lahan
yang cukup jelas. Area air terkonsentrasi pada nilai reflektansi rendah
pada kedua band, sedangkan vegetasi menunjukkan nilai tinggi pada band
NIR namun rendah pada band merah. Area urban berada di antara keduanya
dengan distribusi yang lebih menyebar. Pola diagonal yang terbentuk
menunjukkan hubungan antara kedua band, sekaligus memperlihatkan adanya
overlap pada beberapa titik, terutama pada area urban. Hal ini
menunjukkan adanya mixed pixels akibat resolusi spasial Sentinel-2
sebesar 10 meter.

\#\#LIMITATIONS Terdapat beberapa keterbatasan dalam analisis ini.
Pertama, resolusi spasial Sentinel-2 (10 meter) menyebabkan adanya mixed
pixels, terutama pada area urban yang kompleks. Hal ini menyulitkan
dalam memperoleh sampel yang benar-benar homogen. Kedua, pemilihan
polygon dilakukan secara manual sehingga berpotensi menimbulkan bias
subjektif dalam menentukan area representatif. Selain itu, meskipun data
yang digunakan sudah merupakan Level-2A (Bottom of Atmosphere), faktor
atmosfer dan kondisi lingkungan tetap dapat mempengaruhi nilai
reflektansi.

\#\#FUTURE APPLICATION Teknik yang digunakan dalam praktikum ini
memiliki potensi besar untuk aplikasi di masa depan, seperti pemantauan
pertumbuhan kota, analisis ruang terbuka hijau, serta pengelolaan sumber
daya air. Dengan berkembangnya teknologi seperti Google Earth Engine,
analisis penginderaan jauh dapat dilakukan secara lebih cepat dan dalam
skala yang lebih besar.

Hal ini menunjukkan bahwa penginderaan jauh memiliki peran penting dalam
mendukung pengambilan keputusan berbasis data, khususnya dalam konteks
pembangunan berkelanjutan.

\#\#REFLECTION Pada awal praktikum, penggunaan SNAP terasa cukup
menantang, terutama dalam memahami fungsi berbagai tools seperti spatial
subset, pembuatan polygon, dan analisis spektral. Namun, seiring proses
berlangsung, pemahaman terhadap konsep penginderaan jauh menjadi lebih
jelas, terutama dalam menghubungkan teori dengan hasil visual.

Salah satu hal yang paling menarik adalah melihat bagaimana perbedaan
spektral dapat digunakan untuk membedakan objek secara nyata. Namun,
kesulitan juga muncul pada identifikasi area urban karena sifatnya yang
tidak homogen.

Praktikum ini memberikan pemahaman yang lebih mendalam bahwa
penginderaan jauh tidak hanya bergantung pada data, tetapi juga pada
interpretasi dan pemahaman konsep yang kuat.

\bookmarksetup{startatroot}

\chapter{}\label{section-1}

\#Week2

\# WEEK 2: Landsat 8 Sensor\\
\strut \\
\#\# Summary\\
This week focused on understanding the Landsat 8 sensor and its role in
remote sensing.\\
\strut \\
\#\# Application\\
Landsat 8 data can be used for land cover mapping, vegetation
monitoring, and environmental change detection.\\
\strut \\
\#\# Reflection\\
This week helped me understand how remote sensing sensors collect and
provide data for analysis.

\bookmarksetup{startatroot}

\chapter{}\label{section-2}

\#week3

\bookmarksetup{startatroot}

\chapter{}\label{section-3}

\#\#Week 4-5

\bookmarksetup{startatroot}

\chapter{}\label{section-4}

\#\#week6

\bookmarksetup{startatroot}

\chapter{}\label{section-5}

\#\#week7

\bookmarksetup{startatroot}

\chapter{}\label{section-6}

\#\#week8-9

\bookmarksetup{startatroot}

\chapter*{References}\label{references}
\addcontentsline{toc}{chapter}{References}

\markboth{References}{References}

\phantomsection\label{refs}




\end{document}
